% 图契约：峰值附近充电与峰间电容供电；按本题物理量、参考方向和实际模型绘制。
% 双区对照：实际过程或几何为主体，公式/标注说明其计算作用。
\documentclass[tikz,border=6pt]{standalone}
\usepackage{fontspec,xeCJK,amsmath,amssymb}
\IfFontExistsTF{Times New Roman}{\setmainfont{Times New Roman}}{\setmainfont{TeX Gyre Termes}}
\IfFontExistsTF{Heiti SC}{\setCJKmainfont[AutoFakeBold=1.5]{Heiti SC}}{\setCJKmainfont[AutoFakeBold=1.5]{FandolHei-Regular}}
\usetikzlibrary{arrows.meta,calc}
\definecolor{ink}{HTML}{243745}
\definecolor{blue}{HTML}{275DAB}
\definecolor{green}{HTML}{177C63}
\definecolor{red}{HTML}{BD4935}
\definecolor{orange}{HTML}{AA6C16}
\definecolor{purple}{HTML}{715393}
\begin{document}
\begin{tikzpicture}[x=1cm,y=1cm,font=\fontsize{9.5}{12}\selectfont,
v/.style={-{Latex[length=2mm,width=1.3mm]},line width=.85pt},
axis/.style={-{Latex[length=1.7mm,width=1.1mm]},ink!65,line width=.6pt},
aux/.style={ink!35,dashed,line width=.5pt},
box/.style={draw=ink!35,fill=blue!3,rounded corners=2pt,inner sep=5pt,align=center},
lab/.style={fill=white,inner sep=1.5pt}]
\path[use as bounding box] (-8,-3.6) rectangle(8,3.65);
\node[font=\bfseries\fontsize{10.5}{14}\selectfont,text width=15cm,align=center] at(0,3.23){M6-069-P1　峰值附近充电与峰间电容供电};
\node[text=red] at(-4.25,2.5){正半周峰值附近：充电};
\node[text=green] at(3.55,2.5){两峰之间：桥截止，电容供电};
\begin{scope}[shift={(-4.8,0)}]
\draw[ink](-2,0)--(-1.4328342000348795,0.46082221247166033);
\draw[ink](-0.9671657999651203,0.8391777875283397)--(-0.4,1.3);
\begin{scope}[shift={(-1.2,0.65)},rotate=39.0938588862295]
\draw[ink,fill=white](-.22,-.15)--(-.22,.15)--(.17,0)--cycle;\draw[ink](.17,-.19)--(.17,.19);
\draw[ink](-0.3,0)--(-0.22,0);
\draw[ink](0.17,0)--(0.3,0);
\end{scope}
\node[lab] at(-1.420707418783063,0.921639900040693){$D_1$};
\draw[ink](1.2,0)--(0.6328342000348797,0.46082221247166033);
\draw[ink](0.16716579996512032,0.8391777875283397)--(-0.4,1.3);
\begin{scope}[shift={(0.39999999999999997,0.65)},rotate=140.9061411137705]
\draw[ink,fill=white](-.22,-.15)--(-.22,.15)--(.17,0)--cycle;\draw[ink](.17,-.19)--(.17,.19);
\draw[ink](-0.3,0)--(-0.22,0);
\draw[ink](0.17,0)--(0.3,0);
\end{scope}
\node[lab] at(0.17929258121693697,0.37836009995930714){$D_2$};
\draw[ink](-0.4,-1.3)--(-0.9671657999651203,-0.8391777875283397);
\draw[ink](-1.4328342000348795,-0.46082221247166033)--(-2,0);
\begin{scope}[shift={(-1.2,-0.65)},rotate=140.9061411137705]
\draw[ink,fill=white](-.22,-.15)--(-.22,.15)--(.17,0)--cycle;\draw[ink](.17,-.19)--(.17,.19);
\draw[ink](-0.3,0)--(-0.22,0);
\draw[ink](0.17,0)--(0.3,0);
\end{scope}
\node[lab] at(-1.420707418783063,-0.921639900040693){$D_3$};
\draw[ink](-0.4,-1.3)--(0.16716579996512032,-0.8391777875283397);
\draw[ink](0.6328342000348797,-0.46082221247166033)--(1.2,0);
\begin{scope}[shift={(0.39999999999999997,-0.65)},rotate=39.0938588862295]
\draw[ink,fill=white](-.22,-.15)--(-.22,.15)--(.17,0)--cycle;\draw[ink](.17,-.19)--(.17,.19);
\draw[ink](-0.3,0)--(-0.22,0);
\draw[ink](0.17,0)--(0.3,0);
\end{scope}
\node[lab] at(0.17929258121693697,-0.37836009995930714){$D_4$};
\draw[ink](-2,0)--(-0.73,0);
\draw[ink](-0.07,0)--(1.2,0);
\draw[ink](-.4,0)circle(.33);
\node[] at(-0.4,0){$\sim$};
\draw[ink,line width=.8pt](-0.4,1.3)--(2.9,1.3);
\draw[ink,line width=.8pt](2.9,-1.3)--(-0.4,-1.3);
\draw[ink](1.85,1.3)--(1.85,0.3);
\draw[ink](1.85,-0.3)--(1.85,-1.3);
\begin{scope}[shift={(1.85,0.0)},rotate=-90.0]
\draw[ink](-0.08,-0.25)--(-0.08,0.25);
\draw[ink](0.08,-0.25)--(0.08,0.25);
\draw[ink](-0.3,0)--(-0.08,0);
\draw[ink](0.08,0)--(0.3,0);
\end{scope}
\node[] at(2.2,0.4){$C$};
\draw[ink](2.9,1.3)--(2.9,0.3);
\draw[ink](2.9,-0.3)--(2.9,-1.3);
\begin{scope}[shift={(2.9,0.0)},rotate=-90.0]
\draw[ink,fill=white](-.3,-.12)rectangle(.3,.12);
\end{scope}
\node[lab,anchor=west] at(3.35,0.0){$R_L$};
\node[] at(-2.2,-0.3){$A$};
\node[] at(1.15,-0.3){$B$};
\node[] at(-0.4,1.55){$+$};
\node[] at(-0.4,-1.55){$-$};
\draw[v,green](3.95,0.6)--(3.95,-0.45);
\node[text=green] at(3.95,0.95){$i_L$};
\draw[v,red](-1.8,0.1625)--(-1.45,0.446875);
\draw[v,red](0.78,-0.34125)--(1.1,-0.08125);
\draw[v,red](1.35,0.9)--(1.35,0.2);
\node[text=red] at(1,0.95){$i_C$};
\node[text=red] at(0.4,-2.15){$u_2>u_C+2U_D$};
\end{scope}
\begin{scope}[shift={(3.0,0)}]
\draw[ink!40](-2,0)--(-1.4328342000348795,0.46082221247166033);
\draw[ink!40](-0.9671657999651203,0.8391777875283397)--(-0.4,1.3);
\begin{scope}[shift={(-1.2,0.65)},rotate=39.0938588862295]
\draw[ink!40,fill=white](-.22,-.15)--(-.22,.15)--(.17,0)--cycle;\draw[ink!40](.17,-.19)--(.17,.19);
\draw[ink!40](-0.3,0)--(-0.22,0);
\draw[ink!40](0.17,0)--(0.3,0);
\end{scope}
\node[lab] at(-1.420707418783063,0.921639900040693){$D_1$};
\draw[ink!40](1.2,0)--(0.6328342000348797,0.46082221247166033);
\draw[ink!40](0.16716579996512032,0.8391777875283397)--(-0.4,1.3);
\begin{scope}[shift={(0.39999999999999997,0.65)},rotate=140.9061411137705]
\draw[ink!40,fill=white](-.22,-.15)--(-.22,.15)--(.17,0)--cycle;\draw[ink!40](.17,-.19)--(.17,.19);
\draw[ink!40](-0.3,0)--(-0.22,0);
\draw[ink!40](0.17,0)--(0.3,0);
\end{scope}
\node[lab] at(0.17929258121693697,0.37836009995930714){$D_2$};
\draw[ink!40](-0.4,-1.3)--(-0.9671657999651203,-0.8391777875283397);
\draw[ink!40](-1.4328342000348795,-0.46082221247166033)--(-2,0);
\begin{scope}[shift={(-1.2,-0.65)},rotate=140.9061411137705]
\draw[ink!40,fill=white](-.22,-.15)--(-.22,.15)--(.17,0)--cycle;\draw[ink!40](.17,-.19)--(.17,.19);
\draw[ink!40](-0.3,0)--(-0.22,0);
\draw[ink!40](0.17,0)--(0.3,0);
\end{scope}
\node[lab] at(-1.420707418783063,-0.921639900040693){$D_3$};
\draw[ink!40](-0.4,-1.3)--(0.16716579996512032,-0.8391777875283397);
\draw[ink!40](0.6328342000348797,-0.46082221247166033)--(1.2,0);
\begin{scope}[shift={(0.39999999999999997,-0.65)},rotate=39.0938588862295]
\draw[ink!40,fill=white](-.22,-.15)--(-.22,.15)--(.17,0)--cycle;\draw[ink!40](.17,-.19)--(.17,.19);
\draw[ink!40](-0.3,0)--(-0.22,0);
\draw[ink!40](0.17,0)--(0.3,0);
\end{scope}
\node[lab] at(0.17929258121693697,-0.37836009995930714){$D_4$};
\draw[ink!40](-2,0)--(-0.73,0);
\draw[ink!40](-0.07,0)--(1.2,0);
\draw[ink](-.4,0)circle(.33);
\node[] at(-0.4,0){$\sim$};
\draw[ink,line width=.8pt](-0.4,1.3)--(2.9,1.3);
\draw[ink,line width=.8pt](2.9,-1.3)--(-0.4,-1.3);
\draw[ink](1.85,1.3)--(1.85,0.3);
\draw[ink](1.85,-0.3)--(1.85,-1.3);
\begin{scope}[shift={(1.85,0.0)},rotate=-90.0]
\draw[ink](-0.08,-0.25)--(-0.08,0.25);
\draw[ink](0.08,-0.25)--(0.08,0.25);
\draw[ink](-0.3,0)--(-0.08,0);
\draw[ink](0.08,0)--(0.3,0);
\end{scope}
\node[] at(2.2,0.4){$C$};
\draw[ink](2.9,1.3)--(2.9,0.3);
\draw[ink](2.9,-0.3)--(2.9,-1.3);
\begin{scope}[shift={(2.9,0.0)},rotate=-90.0]
\draw[ink,fill=white](-.3,-.12)rectangle(.3,.12);
\end{scope}
\node[lab,anchor=west] at(3.35,0.0){$R_L$};
\node[] at(-2.2,-0.3){$A$};
\node[] at(1.15,-0.3){$B$};
\node[] at(-0.4,1.55){$+$};
\node[] at(-0.4,-1.55){$-$};
\draw[v,green](3.95,0.6)--(3.95,-0.45);
\node[text=green] at(3.95,0.95){$i_L$};
\draw[v,green](1.35,-0.55)--(1.35,0.25);
\node[text=green] at(0.9,-0.75){$-i_C$};
\node[text=green] at(0.4,-2.15){$|u_2|<u_C+2U_D$};
\end{scope}
\node[text width=15cm,align=center] at(0,-3.12){二极管为恒压降 $U_D$ 模型；负半周峰值换为 $D_2,D_3$ 充电。电容放电期间负载电流仍向下。};
\end{tikzpicture}\end{document}